% 本文件由 LaTeX 排版 Agent 根据有效 translation.json 自动生成。
% author=Caius; work_id=0510; title=卡尤斯残篇

\chapter{卡尤斯残篇}

% structure: p0001 source=h1 target=omitted reason=page-title-already-rendered-by-parent

% structure: p0002 source=h2 target=section reason=relative-heading-rank; base=h2

\section{出自\WorkTitle{驳普罗克路斯对话或论辩}}

% structure: p0003 source=h3 target=subsection reason=relative-heading-rank; base=h2

\subsection{一}

（保存于\Person{欧瑟伯}的\WorkTitle{教会史}，第二卷第二十五章）\par

并且，我可以展示宗徒们的胜利纪念碑。因为如果你愿意前往\Place{梵蒂冈}或\Place{奥斯提亚大道}，你就会找到建立这座教会之人的纪念物。\par

% structure: p0006 source=h3 target=subsection reason=relative-heading-rank; base=h2

\subsection{二}

（同上，第三卷第二十八章）\par

但是，\Person{克林妥}也通过一部据他所说由一位伟大宗徒所写的默示录，向我们呈现奇妙的事，他谎称这些事是由天使向他显示的；他断言在复活之后，基督的王国将在地上，并且\Place{耶路撒冷}中居住的肉体将再次受欲望和享乐的支配。作为天主圣经的仇敌，他企图欺骗人，说将有一千年的时间用于婚宴。\par

% structure: p0009 source=h3 target=subsection reason=relative-heading-rank; base=h2

\subsection{三}

（同上，第三卷第三十一章）\par

此后，在\Place{亚细亚}的\Place{希拉波利斯}，有四位女先知，是\Person{斐理伯}的女儿。她们的坟墓在那里，她们父亲的坟墓也在那里。\par

% structure: p0012 source=h2 target=section reason=relative-heading-rank; base=h2

\section{驳阿特蒙异端}

% structure: p0013 source=h3 target=subsection reason=relative-heading-rank; base=h2

\subsection{一}

（保存于\Person{欧瑟伯}的\WorkTitle{教会史}，第五卷第二十八章）\par

因为他们说，所有那些最初世代的人，甚至宗徒们自己，都领受并教导这些\Emph{人}现在所持守的道理；并且福音宣讲的真理一直保存到\Person{维笃}的时代，他是从\Person{伯多禄}算起的第十三位\Place{罗马}主教，而从他的继任者\Person{则斐利诺}开始，真理就被篡改了。也许他们所说的可能可信，要不是首先圣经本身反驳了他们。此外，还有比\Person{维笃}时代更早的一些弟兄的著作，他们为捍卫真理而反对异教徒，并反驳他们时代的异端：我指的是\Person{犹斯定}、\Person{米提亚德}、\Person{塔提安}、\Person{克莱孟}以及许多其他人，所有这些著作都将天主性归于基督。因为谁不知道\Person{依肋内}、\Person{墨利托}以及其他人的著作，这些著作宣告基督既是天主又是人呢？并且，所有圣咏和赞美诗，这些从一开始由信友们写成的弟兄们的作品，都赞美基督、天主圣言，将天主性归于祂。既然教会教义已经在许多年前被宣告，那么这些人怎么可能像这些\Emph{人}所断言的那样，宣讲到\Person{维笃}的时代呢？他们怎么不羞于对\Person{维笃}进行这些诽谤呢？他们很清楚\Person{维笃}将鞣皮匠\Person{德奥多托}——这个否认天主的背教的领袖和父亲——逐出了教会，因为\Person{德奥多托}首先断言基督只是一个凡人。因为如果像他们所声称的那样，\Person{维笃}持守他们亵渎的教理所教导的同样观点，他怎么会抛弃这个异端的创始人\Person{德奥多托}呢？\par

% structure: p0016 source=h3 target=subsection reason=relative-heading-rank; base=h2

\subsection{二}

（同上，如上）\par

无论如何，我要提醒许多弟兄一件发生在我们自己时代的事——这件事，如果发生在\Place{索多玛}，我想，甚至对他们来说也可能是一个警告。有一位名叫\Person{纳塔利}的宣信者，他不是生在遥远的时代，而是生在我们自己的时代。他曾被\Person{阿斯克勒庇奥多托}和另一个银行家\Person{德奥多托}所欺骗。这两人都是鞣皮匠\Person{德奥多托}的门徒，\Person{德奥多托}是因这种见解（或者更确切地说是愚蠢）而首先被当时的主教\Person{维笃}断绝共融的。\Person{纳塔利}被他们说服，让自己被选为这个异端的主教，条件是他每月从他们那里得到一百五十\OriginalTerm{第纳尔}{denarii}的薪金。于是他与他们联合，并且主多次在异象中劝诫他。因为我们仁慈的天主和主耶稣基督不愿意祂自己苦难的一个见证人因离开教会而丧亡。但他对这些异象不大留意，被在他们中间作首领的尊严以及败坏许多人的卑鄙贪财心所缠绕，最后他在一整夜被圣天使鞭打，严重受击，以至于他清晨起来，披着麻衣，蒙着灰，急急忙忙、泪流满面地投身在主教\Person{则斐利诺}的脚前，不仅跌倒在圣职人员脚前，甚至也跌倒在平信徒脚前，他的哭泣感动了仁慈基督的怜悯教会。在多次恳求并展示他所受击打的伤痕后，他终于被勉强获准共融。\par

% structure: p0019 source=h3 target=subsection reason=relative-heading-rank; base=h2

\subsection{三}

（同上，如上）\par

他们胆敢伪造圣经，弃绝古代信德之准则，忽视基督，不探求圣经所言，反竭力寻求如何构成某种三段论以建立其不敬。若有人向他们提出天主的圣言，他们便考察它是否构成一个连贯或分离的三段论形式；而他们舍弃天主的圣经，钻研几何学，如同属地上的人，谈论地上之事，且不认识那从上而来的。\OriginalTerm{厄乌克肋得斯}{Euclid}实在被他们中一些人费力地丈量，\OriginalTerm{亚里士多德}{Aristotle}和\OriginalTerm{特奥弗拉斯托斯}{Theophrastus}受钦佩，\OriginalTerm{加伦}{Galen}甚至几乎被他们中一些人崇拜。至于那些滥用不信者的技艺以建立自己异端教义，并藉不敬之人的狡猾玷污圣经纯朴信德的人，何须说他们远离信德呢？正因如此，他们胆敢伸手于圣经，声称已加以订正。我并非诬告他们，任何人愿意都可以查证。若有人愿意收集并比较他们所有的抄本，必会发现许多不一致。至少\OriginalTerm{阿斯克勒庇亚德斯}{Asclepiades}的抄本与\OriginalTerm{特奥多图斯}{Theodotus}的抄本相左。许多这样的抄本存在，因为他们的门徒热衷于插入所谓的订正，即他们各自所作的窜改。再者，\OriginalTerm{赫尔默菲洛斯}{Hermophilus}的抄本与这些也不一致；而\OriginalTerm{阿波罗纽斯}{Apollonius}的抄本甚至本身也不一致。因为有人可将他们先前预备的抄本与后来故意窜改的抄本相比较，会发现许多差异。至于这种罪过所含的大胆，他们自己恐怕也不能不知。因为他们要么不信圣经是由圣神默示而成，如此便是不信者；要么自以为比圣神更智慧，那他们岂不是附魔者？他们不能否认这罪行是他们的，因为抄本是出于他们自己的手；他们也没有从最初在信德中教导他们的人那里接受这样的圣经抄本，且不能出示据以抄录的原稿。他们中有些人甚至不屑于窜改，只因他们不法的和不敬的教义，假借恩宠的托辞，干脆否认法律和先知，沉沦到灭亡的最深渊。\par

% structure: p0022 source=h2 target=section reason=relative-heading-rank; base=h2

\section{穆拉托里正典}

（见于\Person{穆拉托里}，\WorkTitle{古意大利中世纪文献}，卷三，第854栏）\par

……他所亲历的那些事他如此记载。福音的第三卷，即路加所记，那位著名的医生\OriginalTerm{路加}{Luke}在基督升天后按顺序以自己的名义写成，且保禄曾视他为追求正义的同伴而与自己结合。他自己并未亲眼见过主在世，而按他所能成就的，从若翰的诞生开始叙述。第四卷福音是\OriginalTerm{若望}{John}的，他是门徒之一。当他的同门和主教们恳求他时，他说：\Quote{你们与我一同禁食三天，让我们彼此述说各人所得之启示。}当夜，宗徒之一的\OriginalTerm{安德肋}{Andrew}得到启示，若望应按他们所记的，以自己的名义叙述一切。因此，尽管福音各书中教导了不同的要点，但在信徒的信德上并无差异，因为所有书中，一切都在同一统御之神下关联着主的诞生、祂的受难、祂的复活、祂与门徒的交谈，以及祂的两次来临——第一次是在被弃绝的卑微中，现已过去；第二次是在王权的荣耀中，尚在未来。那么，若望在书信中也如此不断地提出这些事\BibleRef{若一 1:1}，他亲自说：\BibleQuote{我们亲眼看见过，亲耳听见过，亲手摸过的，都写下来了。}这样，他自称不仅是目击者，也是听者；此外，还是按顺序记录一切有关主之奇事的史家。\par

此外，全部宗徒大事录由路加汇录入一卷，并献给最卓越的\OriginalTerm{德敖斐罗}{Theophilus}，因为这些不同的事件发生在他亲自在场之时；他清楚表明这一原则——即他写作的根据是只记录他所亲见之事——通过略去伯多禄的受难以及保禄从罗马城前往西班牙的旅程。\par

3. 至于\Person{保禄}的书信，对于那些愿意明白此事的人，这些书信本身便指明了自己的性质，以及它们是从何地、为何目的而写的。他首先写了一封相当长的信给\Place{格林多}人，以制止异端的分裂；然后给\Place{迦拉达}人，以禁止割损；再后给\Place{罗马}人，论及\Emph{\WorkTitle{旧约圣经}}的准则，并向他们表明基督是这些经书的首要对象——这需要我们逐一讨论，正如蒙福的宗徒\Person{保禄}，遵循其前任\Person{若望}的准则，只指名写给七个教会，依此顺序：第一封致\Place{格林多}人，第二封致\Place{厄弗所}人，第三封致\Place{斐理伯}人，第四封致\Place{哥罗森}人，第五封致\Place{迦拉达}人，第六封致\Place{得撒洛尼}人，第七封致\Place{罗马}人。此外，虽然他两次写给\Place{格林多}人和\Place{得撒洛尼}人，为纠正他们，但这七封书信却表明——即通过这七重写作——有一个遍布全世界的教会。同样，\Person{若望}在\WorkTitle{默示录}中，虽然只写给七个教会，却向所有人发言。除了这些，他还写了一封给\Person{费肋孟}，一封给\Person{弟铎}，两封给\Person{弟茂德}，这确实是出于单纯个人的情感与爱；但这些书信在天主教会的尊崇中，以及在教会纪律的规范中，也被视为神圣。此外，还有一封致\Place{劳狄刻雅}人的书信，和另一封致\Place{亚历山大里亚}人的书信，伪造\Person{保禄}之名，针对\Person{马西昂}的异端而作；还有其他几封书信不能为天主教会所接纳，因为苦胆与蜂蜜相混是不合适的。\par

4. \Person{犹达}书信，以及前述\Person{若望}的两封书信——或托名\Person{若望}的书信——被列入公函之中。还有由\Person{撒罗满}的朋友们为尊荣他而写的\WorkTitle{智慧书}。我们也接受\Person{若望}的\WorkTitle{默示录}和\Person{伯多禄}的\WorkTitle{默示录}，尽管我们中间有些人不愿后者在教会中宣读。此外，\Person{赫尔玛}在不久前的我们的时代，在\Place{罗马}城，当他的兄弟\Person{普卢斯}主教坐镇罗马教会之椅时，写了\WorkTitle{牧者}。因此，它也应当被诵读；但它不能在教会中公开于民众，也不能列入先知书中，因其数目已满，也不能列入宗徒书中，直到世界的终结。至于\Person{阿尔西诺斯}（又名\Person{瓦伦提努斯}）或\Person{米尔提亚德斯}的著作，我们一概不接受。那些为\Person{马西昂}编写新\WorkTitle{圣咏集}的人，连同\Person{巴西里德}和亚洲卡塔夫里基亚派的创始人，也被拒绝。\par

\SourceRecord{Translated by S.D.F. Salmond；Ante-Nicene Fathers；Vol. 5.；Edited by Alexander Roberts, James Donaldson, and A. Cleveland Coxe.；Buffalo, NY: Christian Literature Publishing Co.,；1886.；<http://www.newadvent.org/fathers/0510.htm>.}
